\documentclass[11pt,a4paper]{article}

% --- Packages ---
\usepackage{amsmath,amssymb,amsthm}
\usepackage{mathtools}
\usepackage[margin=3cm]{geometry}
\usepackage[colorlinks=true,linkcolor=black,citecolor=black,urlcolor=black]{hyperref}
\usepackage{microtype}
\usepackage{booktabs}
\usepackage{array}
\usepackage{enumitem}

% --- Theorem environments (plain LaTeX, no colour boxes) ---
\theoremstyle{plain}
\newtheorem{theorem}{Theorem}[section]
\newtheorem{proposition}[theorem]{Proposition}
\newtheorem{lemma}[theorem]{Lemma}
\newtheorem{corollary}[theorem]{Corollary}

\theoremstyle{definition}
\newtheorem{definition}[theorem]{Definition}
\newtheorem{example}[theorem]{Example}

\theoremstyle{remark}
\newtheorem{remark}[theorem]{Remark}

% --- Macros ---
\newcommand{\ZZ}{\mathbb{Z}}
\newcommand{\Ztwo}{\mathbb{Z}_2}
\newcommand{\Zd}{\mathbb{Z}_d}
\newcommand{\Zhat}{\widehat{\mathbb{Z}}_d}
\newcommand{\Hil}{\mathcal{H}}
\newcommand{\Cft}{\mathcal{C}}
\newcommand{\ii}{\mathrm{i}}
\newcommand{\ee}{\mathrm{e}}
\newcommand{\partfn}[2]{Z\!\begin{bmatrix}#1\\#2\end{bmatrix}}
\newcommand{\Tr}{\operatorname{Tr}}
\newcommand{\abs}[1]{\lvert #1 \rvert}

% --- Title ---
\title{Orbifold Duality in Two-Dimensional CFT\\[4pt]
\large Proof and Explicit Examples}
\author{Based on Vafa (1989) and Dixon--Ginsparg--Harvey (1988)}
\date{}

\begin{document}

\maketitle
\tableofcontents

\bigskip\noindent
\textit{Abstract.}
We prove that for a two-dimensional CFT $\Cft$ with a finite abelian
symmetry $G\cong\Zd$, the orbifold $\Cft/\Zd$ carries a canonical dual
symmetry $\widehat{G}$, and that gauging $\widehat{G}$ recovers $\Cft$.
We verify this in two explicit examples: the free compact boson and the
Ising CFT.

% =====================================================================
\section{Setup}
% =====================================================================

\subsection{Twisted partition functions}

Let $\Cft$ be a two-dimensional CFT on the torus $T^2 = S^1 \times S^1$,
with modular parameter $\tau$ and $q = \ee^{2\pi\ii\tau}$.  For a finite
symmetry group $G$ with elements labelled $h, g$, define the
\emph{twisted partition function}
\begin{equation}\label{eq:twisted}
  \partfn{h}{g}
  \;=\;
  \Tr_{\Hil_h}\!\bigl(g\,q^{L_0 - c/24}\,\bar{q}^{\bar{L}_0 - \bar{c}/24}\bigr),
\end{equation}
where $\Hil_h$ is the Hilbert space with $h$-twisted spatial boundary
condition, and $g$ is inserted as an operator to twist the temporal
boundary condition.

\subsection{The orbifold}

For $G = \Zd = \langle g \rangle$, the \emph{orbifold} partition function is
\begin{equation}\label{eq:orbifold}
  Z_{\Cft/\Zd}
  \;=\;
  \frac{1}{d}\sum_{h,g=0}^{d-1}\partfn{h}{g}.
\end{equation}
This projects onto $G$-invariant states and sums over all twisted
sectors; the result is modular invariant.

% =====================================================================
\section{The Dual Symmetry and the Double-Orbifold Theorem}
% =====================================================================

\subsection{The dual symmetry}

After orbifolding, the Hilbert space decomposes into twisted sectors
\[
  \Hil_{\Cft/\Zd} = \bigoplus_{r=0}^{d-1} \Hil_r^{G\text{-inv}},
\]
where $\Hil_r^{G\text{-inv}}$ is the $G$-invariant subspace of the
$r$-twisted sector.

\begin{definition}\label{def:dual}
  The \emph{dual symmetry generator} $\hat{g}$ acts on
  $\Hil_r^{G\text{-inv}}$ by the phase
  \[
    \hat{g}\,\lvert\psi\rangle_r \;=\; \ee^{2\pi\ii r/d}\,\lvert\psi\rangle_r.
  \]
\end{definition}

\begin{proposition}
  $\hat{g}$ is a well-defined $\Zd$ symmetry of $\Cft/\Zd$.
\end{proposition}
\begin{proof}
  \emph{Order.}  $\hat{g}^d = \ee^{2\pi\ii r} = 1$.
  \emph{Homomorphism.}  The OPE maps the sector-$r_1$ and sector-$r_2$
  subspaces into the sector-$(r_1{+}r_2)$ subspace (string joining
  requires matching boundary conditions), so $\hat{g}$ acts
  multiplicatively on products.
  \emph{Commutativity with $H$.}  $\hat{g}$ assigns a constant phase
  within each energy level, so $[H,\hat{g}]=0$.
\end{proof}

The twisted partition functions of $\Cft/\Zd$ under $\hat{g}$ are
\begin{equation}\label{eq:twisted-orb}
  Z_{\Cft/\Zd}\!\begin{bmatrix}\hat{h}\\\hat{g}\end{bmatrix}
  \;=\;
  \frac{1}{d}\sum_{h,g=0}^{d-1}
  \ee^{2\pi\ii\hat{h}g/d}\;\partfn{h}{g},
\end{equation}
where the phase arises because a spatial $\hat{g}^{\hat{h}}$-twist
multiplies states in sector $g$ by $\ee^{2\pi\ii\hat{h}g/d}$.

\subsection{The double-orbifold theorem}

\begin{theorem}[Double Orbifold]\label{thm:main}
  $\displaystyle Z_{\bigl(\Cft/\Zd\bigr)/\Zhat} \;=\; Z_{\Cft}.$
\end{theorem}

\begin{proof}
Apply~\eqref{eq:orbifold} to $\Cft/\Zd$ and substitute~\eqref{eq:twisted-orb}:
\begin{align}
  Z_{(\Cft/\Zd)/\Zhat}
  &= \frac{1}{d}\sum_{\hat{h},\hat{g}=0}^{d-1}
     Z_{\Cft/\Zd}\!\begin{bmatrix}\hat{h}\\\hat{g}\end{bmatrix}
  = \frac{1}{d^2}
    \sum_{\hat{h},\hat{g},h,g}
    \ee^{2\pi\ii\hat{h}g/d}\;\partfn{h}{g}
  \notag\\[4pt]
  &= \frac{1}{d^2}
    \sum_{\hat{g},h,g}
    \partfn{h}{g}
    \underbrace{\sum_{\hat{h}=0}^{d-1}\ee^{2\pi\ii\hat{h}g/d}}_{d\,\delta_{g,0}}
  = \frac{1}{d}\sum_{\hat{g}=0}^{d-1}\sum_{h=0}^{d-1}
    \partfn{h}{0}
  = \sum_{h=0}^{d-1}\partfn{h}{0}
  = Z_{\Cft}.
  \label{eq:proof}
\end{align}
The key identity is discrete Fourier orthogonality,
\begin{equation}\label{eq:fourier}
  \frac{1}{d}\sum_{\hat{h}=0}^{d-1}\ee^{2\pi\ii\hat{h}g/d} = \delta_{g,\,0\bmod d},
\end{equation}
which kills all $g\neq 0$ sectors.  The remaining sum $\sum_{\hat{g}}$
yields a factor $d$ that cancels the $d^{-2}$ prefactor, and
$\sum_h\partfn{h}{0} = Z_{\Cft}$ because tracing over each twisted
sector without a group insertion reproduces the full untwisted
spectrum.
\end{proof}

\begin{remark}
  The proof is the CFT avatar of Pontryagin duality: $\widehat{\widehat{\Zd}}\cong\Zd$,
  and identity~\eqref{eq:fourier} is Fourier inversion on the group
  algebra $\ZZ[\Zd]$.  For $d=2$ the structure is that of a $2\times 2$
  Hadamard matrix, which squares to $4\cdot\mathbf{1}$.
\end{remark}

% =====================================================================
\section{Example I: The Free Compact Boson}
% =====================================================================

Let $X \sim X + 2\pi R$ with $\alpha' = 1$.  The left/right momenta
for $(m,n)\in\ZZ^2$ are $p_{L,R} = m/R \pm nR/2$, and the partition
function is
\[
  Z_{\Cft}(R)
  = \frac{1}{\abs{\eta(\tau)}^2}
    \sum_{m,n\in\ZZ} q^{\frac{1}{4}p_L^2}\bar{q}^{\frac{1}{4}p_R^2}.
\]

\subsection{The $\Ztwo$ symmetry}

Take $d=2$ with generator $g: X\mapsto X + \pi R$, so
$g\,\lvert m,n\rangle = (-1)^m\lvert m,n\rangle$.

\paragraph{Twisted partition functions.}
\begin{align}
  \partfn{0}{0} &= \frac{1}{\abs{\eta}^2}\sum_{m,n}
    q^{\frac{1}{4}p_L^2}\bar{q}^{\frac{1}{4}p_R^2},
  \label{eq:bos00}\\[4pt]
  \partfn{0}{1} &= \frac{1}{\abs{\eta}^2}\sum_{m,n}
    (-1)^m\,q^{\frac{1}{4}p_L^2}\bar{q}^{\frac{1}{4}p_R^2},
  \label{eq:bos01}\\[4pt]
  \partfn{1}{0} &= \frac{1}{\abs{\eta}^2}
    \left\lvert\vartheta\!\begin{bmatrix}\frac{1}{2}\\0\end{bmatrix}\right\rvert^2
    \sum_{n\in\ZZ}q^{\frac{R^2}{4}(n+\frac{1}{2})^2}
    \bar{q}^{\frac{R^2}{4}(n+\frac{1}{2})^2},
  \label{eq:bos10}\\[4pt]
  \partfn{1}{1} &= \frac{1}{\abs{\eta}^2}
    \left\lvert\vartheta\!\begin{bmatrix}\frac{1}{2}\\\frac{1}{2}\end{bmatrix}\right\rvert^2
    \sum_{n\in\ZZ}(-1)^n\,
    q^{\frac{R^2}{4}(n+\frac{1}{2})^2}
    \bar{q}^{\frac{R^2}{4}(n+\frac{1}{2})^2}.
  \label{eq:bos11}
\end{align}
In~\eqref{eq:bos10}--\eqref{eq:bos11}, the half-integer moding of the
oscillators in the twisted sector is encoded in the Jacobi theta functions
$\vartheta\!\begin{bmatrix}a\\b\end{bmatrix}(\tau)
=\sum_{n\in\ZZ}q^{\frac{1}{2}(n+a)^2}\ee^{2\pi\ii(n+a)b}$,
and the $(-1)^n$ in~\eqref{eq:bos11} comes from the $g$-insertion
acting on winding modes.

\subsection{Orbifold and dual}

The orbifold is $Z_{\Cft/\Ztwo} = \frac{1}{2}\sum_{h,g}\partfn{h}{g}$,
which is equivalent to $Z_{\Cft}(2/R)$ (the $T$-dual theory at radius
$2/R$).  The dual generator $\hat{g}$ acts by $(-1)^r$ on the $r$-twisted
sector.

\subsection{Double-orbifold verification}

Computing the four twisted partition functions of $\Cft/\Ztwo$ under
$\hat{\Ztwo}$ via~\eqref{eq:twisted-orb}:
\begin{align*}
  Z_{\Cft/\Ztwo}\!\begin{bmatrix}0\\0\end{bmatrix}
  &= \tfrac{1}{2}\bigl[\partfn{0}{0}+\partfn{0}{1}+\partfn{1}{0}+\partfn{1}{1}\bigr],\\
  Z_{\Cft/\Ztwo}\!\begin{bmatrix}0\\1\end{bmatrix}
  &= \tfrac{1}{2}\bigl[\partfn{0}{0}+\partfn{0}{1}-\partfn{1}{0}-\partfn{1}{1}\bigr],\\
  Z_{\Cft/\Ztwo}\!\begin{bmatrix}1\\0\end{bmatrix}
  &= \tfrac{1}{2}\bigl[\partfn{0}{0}-\partfn{0}{1}+\partfn{1}{0}-\partfn{1}{1}\bigr],\\
  Z_{\Cft/\Ztwo}\!\begin{bmatrix}1\\1\end{bmatrix}
  &= \tfrac{1}{2}\bigl[\partfn{0}{0}-\partfn{0}{1}-\partfn{1}{0}+\partfn{1}{1}\bigr].
\end{align*}
Summing with weight $\frac{1}{2}$:
\[
  Z_{(\Cft/\Ztwo)/\hat{\Ztwo}}
  = \frac{1}{4}\Bigl[
    4\,\partfn{0}{0}
    + 0\cdot\partfn{0}{1}
    + 0\cdot\partfn{1}{0}
    + 0\cdot\partfn{1}{1}
  \Bigr]
  = \partfn{0}{0} = Z_{\Cft}(R).
\]
The sign cancellations follow directly from the column sums of the
$4\times 4$ Hadamard matrix:
\begin{equation}\label{eq:hadamard}
  \frac{1}{2}\begin{pmatrix}1&1&1&1\\1&1&-1&-1\\1&-1&1&-1\\1&-1&-1&1\end{pmatrix},
  \qquad
  \text{column sums: }2,\;0,\;0,\;0.
\end{equation}
Hence $(\Cft(R)/\Ztwo)/\hat{\Ztwo} \cong \Cft(R)$, consistent with
the $T$-duality chain $\Cft(R)\to\Cft(2/R)\to\Cft(R)$.

% =====================================================================
\section{Example II: The Ising CFT}
% =====================================================================

The Ising CFT is the unitary minimal model $\mathcal{M}(4,3)$ at
$c = \frac{1}{2}$, with three primary fields:
\[
  \mathbf{1}:\;(h,\bar{h})=(0,0),\qquad
  \varepsilon:\;(h,\bar{h})=(\tfrac{1}{2},\tfrac{1}{2}),\qquad
  \sigma:\;(h,\bar{h})=(\tfrac{1}{16},\tfrac{1}{16}).
\]

\subsection{Characters}

Using the Jacobi theta functions
$\vartheta_2 = \vartheta\!\begin{bmatrix}\frac{1}{2}\\0\end{bmatrix}$,
$\vartheta_3 = \vartheta\!\begin{bmatrix}0\\0\end{bmatrix}$,
$\vartheta_4 = \vartheta\!\begin{bmatrix}0\\\frac{1}{2}\end{bmatrix}$,
the Virasoro characters are
\[
  \chi_{\mathbf{1}} = \tfrac{1}{2}\!\left(\sqrt{\tfrac{\vartheta_3}{\eta}}
    +\sqrt{\tfrac{\vartheta_4}{\eta}}\right),\quad
  \chi_{\varepsilon} = \tfrac{1}{2}\!\left(\sqrt{\tfrac{\vartheta_3}{\eta}}
    -\sqrt{\tfrac{\vartheta_4}{\eta}}\right),\quad
  \chi_{\sigma} = \tfrac{1}{\sqrt{2}}\sqrt{\tfrac{\vartheta_2}{\eta}}.
\]
The torus partition function is the diagonal modular invariant:
\begin{equation}\label{eq:ZIsing}
  Z_{\mathrm{Ising}}
  = \abs{\chi_{\mathbf{1}}}^2 + \abs{\chi_\varepsilon}^2 + \abs{\chi_\sigma}^2
  = \partfn{0}{0}.
\end{equation}

\subsection{The $\Ztwo$ symmetry and twisted sectors}

The spin-flip symmetry $g$ acts as
$g: \sigma \mapsto -\sigma$, $g: \mathbf{1},\varepsilon \mapsto
+\mathbf{1},+\varepsilon$.
The four twisted partition functions are:

\medskip
\begin{center}
\renewcommand{\arraystretch}{1.6}
\begin{tabular}{clp{7cm}}
  \toprule
  Sector & Value & Physical origin \\
  \midrule
  $\partfn{0}{0}$ &
  $\abs{\chi_{\mathbf{1}}}^2+\abs{\chi_\varepsilon}^2+\abs{\chi_\sigma}^2$ &
  Untwisted trace, no insertion \\
  $\partfn{0}{1}$ &
  $\abs{\chi_{\mathbf{1}}}^2+\abs{\chi_\varepsilon}^2-\abs{\chi_\sigma}^2$ &
  $g$-eigenvalue of $\sigma$ is $-1$ \\
  $\partfn{1}{0}$ &
  $\abs{\chi_\sigma}^2$ &
  Twisted sector built on disorder operator $\sigma$ \\
  $\partfn{1}{1}$ &
  $-\abs{\chi_\sigma}^2$ &
  All twisted-sector states have $g$-eigenvalue $-1$ \\
  \bottomrule
\end{tabular}
\end{center}

\medskip
One readily checks the modular-transformation identities
$\partfn{h}{g}\to\partfn{h}{g+h}$ under $\tau\to\tau+1$ and
$\partfn{h}{g}\to\partfn{g}{-h}$ under $\tau\to-1/\tau$.

\subsection{The orbifold}

\begin{align}
  Z_{\mathrm{Ising}/\Ztwo}
  &= \tfrac{1}{2}\bigl[
    \partfn{0}{0}+\partfn{0}{1}+\partfn{1}{0}+\partfn{1}{1}
  \bigr]
  \notag\\
  &= \tfrac{1}{2}\bigl[
    2\abs{\chi_{\mathbf{1}}}^2+2\abs{\chi_\varepsilon}^2
    +\abs{\chi_\sigma}^2-\abs{\chi_\sigma}^2
  \bigr]
  = \abs{\chi_{\mathbf{1}}}^2 + \abs{\chi_\varepsilon}^2.
  \label{eq:IsingOrb}
\end{align}
This is the partition function of the \textbf{free Majorana fermion} with
periodic (Neveu--Schwarz) boundary conditions.  The spin field $\sigma$,
being odd under $g$, is projected out.  This orbifold is the
field-theoretic GSO projection.

\subsection{Dual symmetry}

The $g$-twisted Hilbert space of $\mathrm{Ising}$ consists entirely of
states with $g$-eigenvalue $-1$ (since $\sigma$ is odd), so its
$G$-invariant subspace is \emph{empty}.  Consequently
\[
  \Hil_{\mathrm{Ising}/\Ztwo}
  = \Hil_{\mathbf{1}} \oplus \Hil_\varepsilon
  \quad (\text{untwisted sector only, }r=0),
\]
and the dual generator $\hat{g} = (-1)^r = +1$ acts \emph{trivially}
on all states.  Gauging a trivially-acting $\Ztwo$ reintroduces a
new twisted sector without projecting the untwisted one.

\subsection{Double-orbifold verification}

Applying~\eqref{eq:twisted-orb} to compute the four twisted partition
functions of $\mathrm{Ising}/\Ztwo$:
\begin{align}
  Z_{\mathrm{orb}}\!\begin{bmatrix}0\\0\end{bmatrix}
  &= \abs{\chi_{\mathbf{1}}}^2+\abs{\chi_\varepsilon}^2,
  \label{eq:Idd00}\\[3pt]
  Z_{\mathrm{orb}}\!\begin{bmatrix}0\\1\end{bmatrix}
  &= \tfrac{1}{2}\bigl[
    (2\abs{\chi_{\mathbf{1}}}^2+2\abs{\chi_\varepsilon}^2)
    -(-\abs{\chi_\sigma}^2+\abs{\chi_\sigma}^2)
    -(\abs{\chi_\sigma}^2-\abs{\chi_\sigma}^2)
  \bigr]
  = \abs{\chi_{\mathbf{1}}}^2+\abs{\chi_\varepsilon}^2,
  \label{eq:Idd01}\\[3pt]
  Z_{\mathrm{orb}}\!\begin{bmatrix}1\\0\end{bmatrix}
  &= \tfrac{1}{2}\bigl[
    2\abs{\chi_\sigma}^2 + 2\abs{\chi_\sigma}^2
  \bigr] = 2\abs{\chi_\sigma}^2,
  \label{eq:Idd10}\\[3pt]
  Z_{\mathrm{orb}}\!\begin{bmatrix}1\\1\end{bmatrix}
  &= \tfrac{1}{2}\bigl[
    2\abs{\chi_\sigma}^2 - 2\abs{\chi_\sigma}^2
  \bigr] = 0.
  \label{eq:Idd11}
\end{align}
The double orbifold:
\begin{align}
  Z_{(\mathrm{Ising}/\Ztwo)/\hat{\Ztwo}}
  &= \tfrac{1}{2}\bigl[
    (\abs{\chi_{\mathbf{1}}}^2+\abs{\chi_\varepsilon}^2)
    +(\abs{\chi_{\mathbf{1}}}^2+\abs{\chi_\varepsilon}^2)
    +2\abs{\chi_\sigma}^2
    +0
  \bigr]
  \notag\\
  &= \abs{\chi_{\mathbf{1}}}^2+\abs{\chi_\varepsilon}^2+\abs{\chi_\sigma}^2
  = Z_{\mathrm{Ising}}.
  \label{eq:IsingDouble}
\end{align}

The coefficient table makes the cancellation transparent:

\medskip
\begin{center}
\renewcommand{\arraystretch}{1.5}
\begin{tabular}{c|rrrr|l}
  \toprule
  & $\partfn{0}{0}$ & $\partfn{0}{1}$ & $\partfn{1}{0}$ & $\partfn{1}{1}$
  & Result \\
  \midrule
  $Z_{\mathrm{orb}}\!\begin{bmatrix}0\\0\end{bmatrix}$
  & $+1$ & $+1$ & $+1$ & $+1$ & $\abs{\chi_{\mathbf{1}}}^2+\abs{\chi_\varepsilon}^2$ \\
  $Z_{\mathrm{orb}}\!\begin{bmatrix}0\\1\end{bmatrix}$
  & $+1$ & $+1$ & $-1$ & $-1$ & $\abs{\chi_{\mathbf{1}}}^2+\abs{\chi_\varepsilon}^2$ \\
  $Z_{\mathrm{orb}}\!\begin{bmatrix}1\\0\end{bmatrix}$
  & $+1$ & $-1$ & $+1$ & $-1$ & $2\abs{\chi_\sigma}^2$ \\
  $Z_{\mathrm{orb}}\!\begin{bmatrix}1\\1\end{bmatrix}$
  & $+1$ & $-1$ & $-1$ & $+1$ & $0$ \\
  \midrule
  $\frac{1}{2}\sum$ & $2$ & $0$ & $0$ & $0$
  & $Z_{\mathrm{Ising}}$ \\
  \bottomrule
\end{tabular}
\end{center}
\medskip

Only $\partfn{0}{0}$ survives because its column sum is $4$ while all
other column sums vanish -- precisely the discrete Fourier
identity~\eqref{eq:fourier} for $d=2$.

\subsection{Physical interpretation: Jordan--Wigner duality}

The orbifold chain for Ising has a well-known lattice interpretation:
\begin{align*}
  \mathrm{Ising}
  &\;\xrightarrow{\;\text{gauge }g\text{ (spin-flip)}\;}\;
  \text{free Majorana fermion}
  \;\xrightarrow{\;\text{gauge }(-1)^F\;}\;
  \mathrm{Ising}.
\end{align*}
Gauging the spin-flip symmetry is the Jordan--Wigner transformation;
the dual symmetry $\hat{g}=(-1)^F$ is fermion parity, and gauging it
inverts the transformation.

% =====================================================================
\section{Comparison and Summary}
% =====================================================================

\begin{center}
\renewcommand{\arraystretch}{1.6}
\begin{tabular}{l|ll}
  \toprule
  & Compact boson at radius $R$ & Ising CFT \\
  \midrule
  $\Cft$ &
  $Z_{\Cft}(R)$ &
  $\abs{\chi_{\mathbf{1}}}^2+\abs{\chi_\varepsilon}^2+\abs{\chi_\sigma}^2$ \\
  $\Cft/\Ztwo$ &
  $Z_{\Cft}(2/R)$ \;($T$-dual) &
  $\abs{\chi_{\mathbf{1}}}^2+\abs{\chi_\varepsilon}^2$ \;(no $\sigma$) \\
  $\hat{\Ztwo}$ &
  shift $X\mapsto X+\pi\cdot(2/R)$ &
  fermion parity $(-1)^F$ (trivial) \\
  $(\Cft/\Ztwo)/\hat{\Ztwo}$ &
  $Z_{\Cft}(R)$ recovered &
  $Z_{\mathrm{Ising}}$ recovered \\
  Physical name &
  $T$-duality &
  Jordan--Wigner duality \\
  \bottomrule
\end{tabular}
\end{center}

\bigskip
\noindent
The general proof (Theorem~\ref{thm:main}) reduces to the single
identity~\eqref{eq:fourier}: the discrete Fourier transform on $\Zd$
is its own inverse (up to normalisation), reflecting the Pontryagin
self-duality $\widehat{\Zd}\cong\Zd$.  In both examples, only the
untwisted--untwisted sector $\partfn{0}{0}$ survives the double average,
while the three other sectors cancel by destructive interference.

\end{document}
